% ---------------------------------------------------
% Trabajo Final de Grado
% Author: Fabián González Lence <alu0101549491@ull.edu.es>
% Chapter: Modelos de IA y Herramientas
% ----------------------------------------------------

\cleardoublepage
\chapter{Modelos de IA y Herramientas} \label{chap:AIModels}

La elección de los modelos y herramientas de \ac{IA} empleados en este trabajo no es arbitraria: cada uno fue seleccionado en función de sus fortalezas observadas para una fase concreta del ciclo de vida del desarrollo, configurando un equipo de agentes especializados que cooperan de forma coordinada a lo largo del proceso.
Este capítulo describe el panorama actual de los \ac{LLM} aplicados al desarrollo de software, detalla las características y el rol asignado a cada modelo o herramienta utilizada, justifica su selección y documenta el cambio de enfoque metodológico que se introdujo a partir del cuarto proyecto.\\

\section{Panorama actual de los LLMs en el desarrollo de software} \label{sec:llm-overview}

En el momento de realización de este trabajo, el ecosistema de \ac{LLM} orientados al desarrollo de software ha alcanzado un grado de madurez suficiente como para abordar tareas de complejidad real, más allá de la generación de fragmentos de código aislados.
Modelos de propósito general como GPT-4o (OpenAI), Claude 3.5 Sonnet (Anthropic), Gemini 1.5 Pro (Google) o Mistral Large (Mistral AI) compiten directamente en tareas de razonamiento sobre código, diseño arquitectónico y generación de documentación técnica, mientras que herramientas integradas en el entorno de desarrollo como GitHub Copilot han evolucionado desde el autocompletado inteligente hasta un \textbf{modo agente} capaz de ejecutar flujos de trabajo completos de forma autónoma dentro del \ac{IDE}.\\

Esta dualidad entre modelos conversacionales ---accesibles a través de interfaces de chat o API--- y agentes integrados en el entorno de desarrollo resulta especialmente relevante para el presente trabajo, ya que ambas modalidades fueron utilizadas en distintas fases y proyectos.
Los modelos conversacionales ofrecen flexibilidad para el diseño y la supervisión de alto nivel, mientras que los agentes integrados permiten operar directamente sobre el sistema de ficheros, ejecutar comandos, leer errores del compilador y realizar múltiples iteraciones de corrección sin salir del editor.
Esta distinción motivó uno de los hallazgos metodológicos más relevantes del trabajo, que se analiza en detalle en la Sección~\ref{sec:cambio-enfoque}.\\

Cabe señalar también que el panorama de los \ac{LLM} evoluciona a un ritmo muy elevado: los modelos utilizados en este trabajo corresponden a versiones disponibles entre finales de 2024 y principios de 2026, y sus capacidades, limitaciones y modalidades de acceso pueden haber cambiado significativamente en el momento en que el lector consulte esta memoria.
Por este motivo, el análisis que se presenta en este capítulo no busca ofrecer un ranking definitivo de modelos, sino documentar con fidelidad las herramientas concretas que se utilizaron, el rol que se les asignó y el comportamiento que exhibieron en el contexto específico de este experimento.\\

\section{Descripción de los modelos y herramientas de IA} \label{sec:models-description}

A continuación se describen los modelos y herramientas de \ac{IA} empleados a lo largo del trabajo, indicando para cada uno el rol que se le asignó dentro del flujo de desarrollo, las versiones utilizadas y las fortalezas observadas en la práctica.

\subsection{ChatGPT (OpenAI)} \label{subsec:chatgpt}

ChatGPT es el asistente conversacional de OpenAI basado en la familia de modelos GPT-4. Durante el trabajo se utilizó principalmente en las fases iniciales del proyecto como herramienta de apoyo para la redacción de especificaciones de requisitos y la exploración de alternativas arquitectónicas, aprovechando su capacidad para procesar descripciones en lenguaje natural y devolver documentos estructurados.
Su fortaleza más destacada es la fluidez en la generación de texto técnico y su familiaridad con convenciones ampliamente documentadas en la literatura del software.
Sin embargo, en tareas de codificación extensa se observó una tendencia a perder coherencia con el contexto arquitectónico previo cuando la ventana de conversación se alargaba, lo que limitó su utilidad en los proyectos de mayor complejidad.

\subsection{Claude (Anthropic)} \label{subsec:claude}

Claude es el modelo conversacional de Anthropic, y fue el designado como \textbf{\ac{IA} Arquitecto} en el flujo de desarrollo de los cinco proyectos.
En el momento de realización de este trabajo se utilizó principalmente la versión Claude 3.5 Sonnet, y posteriormente Claude Sonnet 4.5, disponible a través de la interfaz web de \textit{claude.ai} y como agente integrado en GitHub Copilot.
Su elección para el rol arquitectónico se fundamentó en tres fortalezas claramente observadas: una \textbf{ventana de contexto amplia} que le permitía procesar diagramas de clases completos y especificaciones de requisitos extensas en una única interacción; una \textbf{notable capacidad para generar documentación técnica detallada} ---incluyendo estructuras de directorios, esqueletos de clases y archivos de configuración--- a partir de descripciones en lenguaje natural; y una \textbf{coherencia arquitectónica} que se mantenía consistente a lo largo de conversaciones largas.\\

La tarea concreta asignada a Claude en cada proyecto consistía en, partiendo de las especificaciones de requisitos y los diagramas de clases y casos de uso generados previamente, generar la estructura completa del proyecto: árbol de directorios, ficheros de configuración (package.json, tsconfig.json, vite.config.ts, eslint.config.mjs, jest.config.js, typedoc.json), esqueletos de clases con sus métodos declarados sin implementar, y documentación base (README, ARCHITECTURE).
El prompt de arquitectura desarrollado para esta fase y detallado en el Capítulo~\ref{chap:Metodologia} incluye instrucciones precisas sobre la pila tecnológica, los patrones de diseño a aplicar y el formato esperado de los artefactos de salida.

\subsection{Mistral} \label{subsec:mistral}

Mistral AI es una empresa francesa fundada en 2023 cuyos modelos, como Mistral Large y Mistral Medium, destacan por su eficiencia y su accesibilidad a través de una interfaz conversacional web.
Durante los tres primeros proyectos (\ac{HANGMAN}, \ac{PLAYER} y \ac{BALATRO}), Mistral fue el designado como \textbf{\ac{IA} Desarrollador}, responsable de implementar el código fuente de cada módulo a partir de los artefactos de diseño producidos por Claude.
Su elección inicial respondía a su buena relación entre capacidad de generación de código y accesibilidad desde una interfaz chatbot sin necesidad de un entorno de desarrollo integrado.\\

En la práctica, Mistral demostró ser capaz de producir implementaciones funcionales para módulos de complejidad media siguiendo fielmente los diagramas de clases proporcionados.
No obstante, a medida que los proyectos escalaron en complejidad, se hicieron evidentes sus limitaciones: la gestión del estado entre módulos interdependientes resultaba inconsistente cuando el contexto superaba ciertos límites, y la generación de código para patrones arquitectónicos avanzados ---como la Clean Architecture con inyección de dependencias--- requería iteraciones numerosas y correcciones manuales significativas.
Estas limitaciones motivaron el cambio de enfoque metodológico descrito en la Sección~\ref{sec:cambio-enfoque}.

\subsection{GitHub Copilot} \label{subsec:copilot}

GitHub Copilot es la herramienta de asistencia al código de GitHub, desarrollada en colaboración con OpenAI, que evoluciona desde un autocompletado inteligente hasta un \textbf{modo agente} capaz de ejecutar tareas complejas de forma autónoma dentro del \ac{IDE}.
A partir del cuarto proyecto (\ac{CARTO}), GitHub Copilot en modo agente pasó a asumir los roles de \textbf{\ac{IA} Desarrollador}, \textbf{\ac{IA} Revisor} y \textbf{\ac{IA} Tester}, sustituyendo a Mistral y Qwen en esas fases.\\

La versión de Copilot utilizada en las fases de codificación operaba con los modelos Claude Sonnet 4.5 y, en algunos momentos, con el modo automático de selección de modelo.
Su principal fortaleza frente al enfoque conversacional previo reside en la \textbf{integración directa con el sistema de ficheros del proyecto}: el agente puede leer el código existente, ejecutar comandos de compilación y linting, interpretar los errores producidos y aplicar correcciones de forma iterativa sin salir del editor.
Esta capacidad eliminó la fricción de copiar y pegar código entre la interfaz chatbot y el \ac{IDE}, y permitió gestionar proyectos de cientos de ficheros con una coherencia que los modelos conversacionales no podían mantener.\\

En el rol de Revisor, Copilot aplicó los prompts estructurados de revisión de código ---desarrollados y documentados en la Guía de Codificación del cuaderno Notion del proyecto--- para generar informes detallados con incidencias clasificadas por severidad y puntuaciones ponderadas.
En el rol de Tester, generó suites de pruebas completas con Jest (proyectos 1--3) y Playwright (proyectos 4--5), incluyendo la gestión del ciclo completo de planificación de casos de prueba, implementación y corrección iterativa de fallos.

\subsection{Perplexity} \label{subsec:perplexity}
Perplexity es un motor de búsqueda potenciado por \ac{IA} que combina la búsqueda en tiempo real con la generación de respuestas referenciadas. Fue evaluado durante las fases iniciales del proyecto como herramienta de consulta; sin embargo, tras las primeras pruebas exploratorias, no llegó a integrarse en el flujo de trabajo del \ac{TFG}. Su uso quedó limitado a consultas aisladas y no aportó ventajas respecto a las herramientas adoptadas, por lo que fue descartado.\\

\subsection{Gemini --- Generación de Imágenes}
\label{subsec:gemini-imagenes}
Google Gemini fue utilizado de forma puntual durante el desarrollo del proyecto \textit{Mini Balatro} (véase la Sección~\ref{sec:balatro}) con un propósito concreto: la \textbf{generación de los assets gráficos} de las cartas de la baraja. Dado que el juego requería un conjunto de imágenes visualmente coherentes para representar las particularidades de las cartas, se recurrió a las capacidades de generación de imagen de Gemini para producir dichos recursos de forma rápida y consistente.\\

Este uso, aunque acotado a un único proyecto y a una fase concreta del desarrollo (la producción de recursos visuales), resultó más significativo en términos prácticos que otras herramientas exploradas durante el \ac{TFG}, al resolver directamente una necesidad real del proyecto sin requerir activos externos ni trabajo manual de diseño.\\

\subsection{Qwen (Alibaba)} \label{subsec:qwen}

Qwen es la familia de modelos de lenguaje de gran tamaño desarrollada por Alibaba Cloud.
En los tres primeros proyectos, Qwen fue el designado como \textbf{\ac{IA} Tester}, responsable de generar las suites de pruebas unitarias con Jest a partir del código revisado y aprobado en la fase anterior.
Su elección para este rol respondía a una observación empírica: Qwen mostraba una particular eficacia en la generación de casos de prueba exhaustivos siguiendo el patrón AAA (\textit{Arrange-Act-Assert}), incluyendo casos normales, límite y excepcionales, con una tasa de compilación inicial elevada.\\

Al igual que Mistral, Qwen fue desplazado a partir del proyecto \ac{CARTO} por GitHub Copilot en modo agente, utilizando modelos como Claude Sonnet 4.5, Claude Sonnet 4.6, GPT-5.2 y GPT-5.4, fundamentalmente porque la complejidad de los proyectos de mayor escala hacía impracticable el flujo de copiar código entre interfaces chatbot para generar pruebas módulo a módulo.
Su contribución en los primeros proyectos fue significativa: para \ac{HANGMAN}, por ejemplo, Qwen generó ficheros de test para los nueve módulos del proyecto, con iteraciones de corrección cuando los tests fallaban, produciendo suites que cubrían los principales casos de uso de cada componente.

\section{Justificación de la selección de herramientas} \label{sec:reason-for-selection}

La configuración del equipo de herramientas de \ac{IA} descrita en las secciones anteriores no fue el resultado de una elección única y definitiva al inicio del trabajo, sino de un proceso iterativo en el que la experiencia acumulada en cada proyecto fue ajustando qué herramienta resultaba más adecuada para cada fase y para cada nivel de complejidad.
A continuación se justifica la asignación de roles adoptada y se documenta el cambio de enfoque metodológico que se introdujo a partir del cuarto proyecto, que constituye uno de los hallazgos más relevantes del trabajo.\\

\subsection{Criterios de asignación de roles} \label{subsec:asignation-criteria}

La asignación de cada herramienta a un rol específico dentro del flujo de cuatro fases ---Arquitectura, Codificación, Revisión y Testing--- respondió a un criterio principal: emplear cada modelo en aquella tarea en que su fortaleza diferencial fuera más aprovechable y su debilidad tuviera menor impacto.\\

Para el rol de \textbf{\ac{IA} Arquitecto}, la elección de Claude se fundamentó en su capacidad para procesar artefactos de diseño extensos ---especificaciones de requisitos, diagramas de clases en formato Mermaid, diagramas de casos de uso--- en una única interacción, y en producir a partir de ellos estructuras de proyecto coherentes con el diseño proporcionado. Esta tarea requería razonamiento arquitectónico sostenido y alta fidelidad al contexto de entrada, cualidades en las que Claude demostró ser consistentemente superior al resto de modelos evaluados durante el trabajo.\\

Para el rol de \textbf{\ac{IA} Desarrollador}, la elección inicial de Mistral respondía a criterios de accesibilidad y a las características del flujo conversacional adoptado para los primeros proyectos: dado que la codificación se realizaba módulo a módulo mediante un chatbot web, se necesitaba un modelo con buena capacidad de generación de código TypeScript, comprensión de diagramas de clases y disposición a seguir instrucciones estructuradas. Mistral cumplía estas condiciones para proyectos de complejidad baja y media, con la ventaja adicional de ser accesible de forma gratuita a través de su plataforma web.\\

Para el rol de \textbf{\ac{IA} Tester}, Qwen fue seleccionado en los primeros proyectos por su eficacia empíricamente observada en la generación de suites de pruebas unitarias siguiendo el patrón \textbf{\ac{AAA}}, con una elevada tasa de compilación inicial. Este rol exigía principalmente capacidad para inferir el comportamiento esperado de cada módulo a partir de su código, generar casos normales, límite y excepcionales, y corregir los tests que fallaban en iteraciones sucesivas; tareas en las que Qwen resultó adecuado para los proyectos de menor escala. Para los últimos dos proyectos, en donde la interacción con el código suponía un problema, se decidió crear un agente de GitHub Copilot\\

El rol de \textbf{\ac{IA} Revisor} fue asignado desde el principio a GitHub Copilot operando con el modelo GPT-4o mini ---siendo cambiado a posteriori por GPT-5.2 y GPT-5.4 en los últimos proyectos---, porque el proceso de revisión requería no solo evaluar código, sino también generar informes estructurados con puntuaciones ponderadas según los criterios de diseño, calidad, cumplimiento de requisitos, mantenibilidad y buenas prácticas ---descritos en detalle en el Capítulo~\ref{chap:Metodology}---. La integración de Copilot con el entorno de desarrollo permitía además que el revisor accediera directamente a los ficheros del proyecto sin necesidad de copiar y pegar código en una interfaz externa.\\

\subsection{El cambio de enfoque metodológico} \label{sec:approach-change}

La transición de los proyectos 1--3 al proyecto 4 (\ac{CARTO}) supuso un punto de inflexión en la metodología del trabajo, motivado por las limitaciones de escalabilidad del flujo conversacional hasta entonces empleado. En los tres primeros proyectos, el agente \ac{IA} Desarrollador (Mistral) operaba como chatbot web: el ingeniero copiaba manualmente el prompt de codificación en la interfaz, obtenía el código generado, lo transfería al editor y repetía el proceso módulo a módulo. Este flujo resultaba viable para proyectos de 9 a 15 módulos con arquitectura relativamente plana; sin embargo, a medida que la complejidad arquitectónica y el número de ficheros crecía ---\ac{CARTO} llegó a gestionar más de trescientos ficheros distribuidos en cinco capas de una arquitectura limpia---, el enfoque comenzó a mostrar limitaciones críticas.\\

Las principales dificultades observadas fueron tres. En primer lugar, la \textbf{pérdida de coherencia inter-módulo}: al generar cada fichero en conversaciones independientes, los modelos conversacionales no podían mantener consistencia entre los contratos de interfaces, los tipos de datos y las convenciones de nomenclatura definidas en módulos anteriores, lo que introducía errores de integración que requerían corrección manual extensiva. En segundo lugar, la \textbf{fricción operativa del flujo manual}: copiar código entre el chatbot y el editor, aplicar cambios, verificar errores de compilación y volver al chatbot con el mensaje de error consumía un tiempo desproporcionado respecto al tiempo de generación real. En tercer lugar, la \textbf{imposibilidad práctica de revisar y testear a esa escala}: intentar revisar más de trescientos ficheros pasando código entre interfaces chatbot habría resultado inviable, tanto en tiempo como en coherencia de los informes resultantes.\\

Ante estas limitaciones, se tomó la decisión de adoptar \textbf{GitHub Copilot en modo agente} como herramienta principal para las fases de Codificación, Revisión y Testing a partir de \ac{CARTO}. El modo agente de Copilot opera directamente sobre el sistema de ficheros del proyecto: puede leer la base de código existente, ejecutar comandos de compilación y linting, interpretar los errores producidos, realizar búsquedas sobre el código fuente con expresiones regulares y aplicar correcciones iterativas sin salir del editor. Esta capacidad de actuar sobre el estado real del proyecto, en lugar de operar sobre representaciones textuales del mismo, resolvió los tres problemas identificados.\\

Durante el periodo de Prácticas Externas, se realizó un curso de introducción a las herramientas que GitHub Copilot tiene disponible, entre las cuales se encontraba \textbf{la personalización de agentes de GitHub Copilot}. Esta herramienta, además de permitir a la \ac{IA} tener gran control sobre el sistema de ficheros del proyecto como lo haría el modo agente normal, permite definir reglas, instrucciones y modelos a utilizar para varios agentes distintos en un formato de fichero Markdown, organizados en un directorio específico dentro del repositorio del proyecto \cite{URL::Agents}. Esto fue de gran utilidad, permitiendo prescindir de los chatbots y pudiendo \textbf{asignar cada uno de los 4 roles principales a un agente personalizado de GitHub Copilot}.\\

Desde un punto de vista práctico, el cambio implicó rediseñar los prompts de las fases de Revisión y Testing para adaptarlos al modelo agente. El \textbf{Revisor 2.0}, aplicado en \ac{CARTO}, adoptó un flujo de tres etapas: en la primera, el agente exploraba toda la base de código y generaba una TODO List con los grupos de ficheros a revisar; en la segunda, revisaba sistemáticamente cada grupo y generaba un informe de incidencias con identificadores, severidades y propuestas de corrección; en la tercera, seguía el informe para implementar y verificar las correcciones. En \ac{CARTO}, este proceso identificó 225 incidencias distribuidas en 41 grupos de revisión, con 5 incidencias críticas, 30 de severidad alta, 111 media y 79 baja. El \textbf{Tester 2.0} adoptó una lógica análoga para la generación de pruebas \textit{end-to-end} con Playwright: el agente analizaba primero los diagramas de casos de uso para definir los escenarios de prueba, y en una segunda pasada implementaba los tests correspondientes; en \ac{CARTO} se generaron 132 escenarios y se implementaron 131 tests reales.\\

Una característica observada como efecto secundario del cambio de herramienta fue que las respuestas de los agentes de Copilot no producen salidas en formato Markdown estándar, sino trazas de ejecución que incluyen referencias a ficheros, resultados de búsquedas y fragmentos de código intercalados. Para paliar esto, se optó por pedirle al propio agente que generase informes de desarrollo resumidos en un formato más legible, o bien que otra \ac{IA} procesara la respuesta bruta para producir un resumen estructurado adecuado para el cuaderno de seguimiento.\\

\section{Comparativa de modelos} \label{sec:comparativa-modelos}

La Tabla~\ref{tab:comparativa-modelos} recoge de forma sintética los modelos y herramientas de \ac{IA} empleados en el trabajo, agrupando para cada uno el rol asignado, las fases del ciclo de desarrollo en que participó, los proyectos en que fue utilizado, sus principales fortalezas, sus limitaciones más relevantes y la modalidad de acceso empleada.
Esta visión de conjunto complementa las descripciones individuales de la sección anterior y sirve de referencia para el análisis de resultados del Capítulo~\ref{chap:Results}.\\

\begin{table}[H]
\centering
\small
\renewcommand{\arraystretch}{1.3}
\begin{tabular}{|p{2.2cm}|p{3.2cm}|p{1.8cm}|p{3cm}|p{3cm}|}
\hline
\textbf{Modelo} & \textbf{Rol y Fases} & \textbf{Proyectos} & \textbf{Fortalezas} & \textbf{Limitaciones} \\
\hline
\textbf{Claude} &
Arquitecto (arquitectura); agente de codificación (P5) &
Todos &
Contexto amplio; coherencia arquitectónica; alta fidelidad al diseño; documentación técnica detallada &
Suscripción de pago; integración como agente dependiente de Copilot \\
\hline
\textbf{Mistral} &
Desarrollador (codificación) &
P1--P3 &
Gratuito; buen rendimiento en módulos de complejidad media; sigue instrucciones estructuradas &
Pérdida de coherencia inter-módulo; dificultad con patrones avanzados; flujo manual no escalable \\
\hline
\textbf{GitHub Copilot} &
Revisor (todos); desarrollador y tester (P4--P5) &
Todos &
Integración directa con ficheros; ejecución iterativa de comandos; gestión coherente a gran escala &
Dependencia de modelos externos; trazas no legibles en Markdown \\
\hline
\textbf{Qwen} &
Tester (testing) &
P1--P3 &
Eficaz con patrón \textit{Arrange-Act-Assert}; alta compilación inicial; cubre casos normales, límite y excepcionales &
Flujo manual no escalable; desplazado por Copilot agente a partir de P4 \\
\hline
\textbf{ChatGPT} &
Apoyo en especificación de requisitos (diseño previo a codificación) &
P1--P3 (en pequeños arreglos de diseño) &
Fluidez en texto técnico; familiaridad con convenciones estándar &
Pérdida de coherencia en conversaciones largas; utilidad decreciente en proyectos complejos \\
\hline
\textbf{Perplexity} &
Consulta técnica puntual (exploratoria) &
Ninguno &
Búsqueda en tiempo real con referencias verificables &
No integrado en el flujo formal; descartado tras pruebas iniciales \\
\hline
\end{tabular}
\caption{Comparativa de modelos y herramientas de \ac{IA} empleados. P1--P5: \ac{HANGMAN}, \ac{PLAYER}, \ac{BALATRO}, \ac{CARTO} y \ac{TENNIS} respectivamente.}
\label{tab:comparativa-modelos}
\end{table}

Varios patrones emergen de esta comparativa.
En primer lugar, ningún modelo resultó adecuado para todas las fases: las fortalezas de Claude en razonamiento arquitectónico lo hacen menos adecuado para la ejecución iterativa de tareas de codificación masiva, mientras que la capacidad operativa de Copilot como agente no sustituye la visión de alto nivel que Claude aporta en la fase de diseño.
En segundo lugar, la modalidad de acceso condiciona fuertemente la elección: la accesibilidad gratuita de Mistral y Qwen resultó suficiente para los proyectos de menor complejidad, pero la escalabilidad del flujo requirió en los proyectos más complejos adoptar herramientas con integración nativa en el entorno de desarrollo.
En tercer lugar, el rol de Revisor fue el único que se mantuvo estable ---siempre asignado a GitHub Copilot--- a lo largo de todos los proyectos, lo que refleja que la integración con el sistema de ficheros es especialmente crítica en las tareas de revisión, donde el acceso directo al código es imprescindible para producir informes de calidad.\\